\documentclass[11pt]{article}
\usepackage{amsmath,amssymb,amsthm,hyperref,geometry}\geometry{margin=1in}
\newtheorem{theorem}{Theorem}
\title{Transfer from an occurrence lift to display-exact content}
\date{15 November 2026}\begin{document}\maketitle
\paragraph{Status.} Proposed relative transfer theorem, not a new independent construction of an occurrence lift. A lift already contains the context's interaction with monadic occurrence content. The theorem below transfers that existing structure to a potentially much richer independent contentad; it must not be advertised as deriving all interaction from an arbitrary contextad.

\section{Exact independent and inherited data}
Retain the base $\mathcal C$, display class $\mathcal D$, Cartesian monad $S$, and independently valid display-exact contentad $(L,p,d,\nu)$ of the display-exact criterion. Replace actual display contexts by a contextad $q:K\to\mathcal C$ with extension $e:K\to\mathcal C$ and comprehension $\epsilon_P:eP\to qP$ in $\mathcal D$. Require comprehension to preserve Cartesian reindexing and to identify context unit and dependent accumulation with identity and composite display maps, through specified coherent isomorphisms. It need not be faithful: different native grade arrows can have the same carrier map.

Separately supply a lift $\overline S:K\to K$ over $S$ with natural comparison $\theta:e\overline S\cong Se$. Require explicitly $\epsilon_{\overline SP}=S(\epsilon_P)\theta_P$, so the comparison is over $SA$, not merely an isomorphism of unrelated carriers. All unit/accumulation comparisons below also commute with their stated comprehension/extension maps. It is a Cartesian functor for $q$; its lifted monad unit and multiplication are Cartesian arrows over $\eta^S$ and $\mu^S$, with extension maps identified by $\theta$ with $\eta^S_e$ and $\mu^S_e$. It preserves context unit and dependent accumulation through coherent native isomorphisms. In particular for $P\in K_A$, $R\in K_{eP}$, its accumulation comparison identifies $\overline S(P\otimes R)$ with $\overline SP\otimes\overline SR$, transporting the latter factor along $\theta_P:e\overline SP\cong SeP$ as necessary.

\emph{Coherence is not omitted:} these comparisons must form the existing native monadic-occurrence interaction, with its full unit, multiplication, naturality and distributive equations. Equivalently one can supply the already proved presented monadic crossing together with its Cartesian lift description and $\theta$. This is precisely the type of existing presentation used in the varying-diagram monadic theorem. A bare functor over $S$, even with $\theta$, is insufficient. No arbitrary content grade $Q\in L$ enters this input, but the monadic $S$ interaction has already been solved.

\section{Transferred crossing}
For $P\in K_A$ and $Q\in L_A$, define
\[
 \widetilde P_Q=\nu_Q^*\overline SP\in K_{dQ},\qquad
 \widetilde Q_P=\epsilon_P^*Q\in L_{eP}.
\tag{T}
\]
The extension of the former is canonically
\[
 e\widetilde P_Q\cong dQ\times_{SA}SeP
                  \cong d(\epsilon_P^*Q).
\tag{B}
\]
The first comparison is Cartesian comprehension plus $\theta$, and the second is the independent content display. They supply the actual body, not merely an endpoint isomorphism class. Representatives with initial body $f:B\to dQ$ and final pure map $g:eP\to C$ use $f^*\widetilde P_Q$, $g_!\widetilde Q_P$, and the composite body from (B).

\begin{theorem}[Relative content transfer]
These data determine a full native crossing for $K$ and $L$, retaining their noninvertible native arrows and pure endpoint squares. Its four comparisons and full coherence are inherited from the existing monadic occurrence law and the independent content Cartesian/display/BC conditions. It is a law on the literal comma family, not on a quotient by extension maps.
\end{theorem}
\begin{proof}
A total native pair consisting of a context arrow $a:P\to P'$ over $u:A\to A'$ and a content arrow $b:Q\to Q'$ over $u$ gives a context map by applying $\overline S$ to $a$ and factoring through the Cartesian pullback along $\nu_{Q'}$. Naturality $Sp(b)\nu_Q=\nu_{Q'}db$ provides the base square. Cartesian uniqueness proves composition. The content map is the Cartesian factorization of $bc_{\epsilon_P,Q}$ through $c_{\epsilon_{P'},Q'}$ over $ea$. The two projections in (B) prove its body square. Pure-middle balance uses Cartesian functoriality of $\overline S$ and the independent canonical content BC mate. This proves descent on all represented generators.

The context unit and context-accumulation comparisons are pullbacks of their existing lifted $S$ comparisons, followed by reindexing compositors. Their content counterparts are composition of Cartesian lifts. The content-unit comparison uses the Cartesian lifted unit of $\overline S$ and the equation $\nu_U\eta^L=\eta^S$; its grade comparison is Cartesian preservation by $U$. The content-accumulation comparison uses Cartesian lifted multiplication of $\overline S$, the equation $\nu_{Q\diamond R}\mu^L=\mu^S S\nu_Q\nu_R$, and Cartesian preservation by content accumulation. These have the same carrier diagrams as in the display-exact criterion, but their native context components come from the supplied lift, not from carrier uniqueness.

For every distributive coherence diagram, compose the context routes with their defining Cartesian arrows into the same lifted monadic expression. The routes agree there by the supplied native monadic-occurrence coherence and reindexing coherence. Cartesian uniqueness then proves equality \emph{in $K$}. In the content component, compose with the target Cartesian arrow and use the independently valid content laws and canonical mate pasting; equality follows \emph{in $L$}. The two comparisons (B) identify the bodies. This proves all ten equations without assuming $e$ or $d$ faithful. Retained comma normalization, including its action on both native arrows and pure balance, is unchanged from the general presented construction.
\end{proof}

\section{Constructive instantiations and what they do not prove}
For the faithful display contextad, $\overline S$ sends $\pi:E\to A$ to $S\pi:SE\to SA$. Pullback preservation and Cartesian $\eta^S,\mu^S$ give the lift; composition of arrows supplies the accumulation comparison. Its coherence is the monad law and pullback uniqueness. Thus the direct display theorem constructs this input rather than assumes it.

For the existing Set rig contexts and finitary addressed polynomial $S$, at a polynomial label $(c,(a_i)_{i\in E_c})$, the lifted context is the addressed tensor of the original context objects $P_{a_i}$. Its extension is $S(eP)$. The monoidal unit treats zero positions, strong extension identifies tensor carriers, dependent substitution uses the rig's canonical finite distributivity, and polynomial unit/multiplication identify unchanged or flattened occurrence addresses. Native maps act at each address; no faithfulness of coefficient extension is required. The original O1 construction proves these native coefficient diagrams, not merely their carrier images. Transfer repackages that existing construction and does not establish a new theorem about arbitrary rigs lacking those assumptions.

For the existing effect-fibre contexts $(X_a,t_a\in TX_a)$ with commutative Set monad $T$, the lift at the same label has raw carrier $\prod_iX_{a_i}$ and effect $\mathsf{dst}(t_{a_i})$. Reindexing keeps occurrence addresses; unit uses a singleton or the original fibre, and flattening identifies nested tuples with leaf tuples. Context accumulation uses the bijection between a product of dependent sums and the dependent sum of products, whose effect equation is commutative bind/Fubini interchange. Those equations, including zero positions and all native maps, were proved in the effect-fibre occurrence manuscript. Normalized disjoint sums are not used, and affinity is not needed for finite crossing. Again the lift is constructively obtained before transferring to arbitrary display-exact content.

These instances verify that the transfer is not empty and can retain nonfaithful grade information. They do not constitute a new arbitrary-dependence theorem: an independently chosen concurrent coupling is not thereby forced to be the addressed product effect. The transfer's value is a clean reusable separation between occurrence interaction and richer content; the old O1 and effect proofs already instantiate much of it. No novelty or full-feedback claim is made.
\end{document}
